% !TEX root = ../main.tex

%----------------------------------------------------------------------------------------
% CHAPTER 10 - IMPLANTACIÓ I RESULTATS
%----------------------------------------------------------------------------------------

\chapter{Implantació i resultats}
\label{Ch:Deployment}

%----------------------------------------------------------------------------------------

\section{Desplegament del sistema}

El sistema s'ha desplegat utilitzant infraestructura gratuïta proporcionada per GitHub, aprofitant que el projecte és de codi obert i està allotjat en aquesta plataforma. L'arquitectura de desplegament es divideix en dues parts:

\subsection{Frontend: GitHub Pages}

L'aplicació frontend s'ha desplegat a GitHub Pages, un servei d'allotjament de llocs web estàtics gratuït proporcionat per GitHub. El desplegament es realitza automàticament mitjançant el workflow \texttt{deploy-frontend-pages.yml}, que:

\begin{enumerate}
    \item Detecta canvis a la branca \texttt{main} o s'executa manualment.
    \item Construeix l'aplicació React amb Vite, generant fitxers estàtics optimitzats al directori \texttt{dist/}.
    \item Puja els fitxers a GitHub Pages mitjançant \texttt{actions/deploy-pages@v4}.
\end{enumerate}

L'aplicació és accessible a l'URL: \url{https://alexruizmaso.github.io/opendata-water-visualization}

Com que es tracta d'una aplicació completament estàtica, no requereix un servidor de backend en producció. Totes les dades es serveixen com a fitxers JSON estàtics des del directori \texttt{public/data/}, eliminant la necessitat de bases de dades o servidors d'aplicacions.

\subsection{ETL: GitHub Actions}

El pipeline ETL s'executa a la infraestructura de GitHub Actions, aprofitant els 2.000 minuts mensuals gratuïts disponibles per a repositoris públics. L'execució es programa diàriament a les 02:00 UTC mitjançant \texttt{cron: '0 2 * * *'}, hora seleccionada per minimitzar l'impacte sobre possibles usuaris i aprofitar les finestres de baixa activitat de les APIs públiques.

Cada execució de l'ETL actualitza els fitxers JSON amb les dades més recents i els desa al repositori, desencadenant automàticament un nou desplegament del frontend si cal.

\subsection{Configuració del domini}

El desplegament utilitza el domini per defecte de GitHub Pages (\texttt{github.io}). No s'ha configurat un domini personalitzat, ja que el domini per defecte és suficient per a un projecte acadèmic i evita costos addicionals. La connexió es realitza automàticament via HTTPS, proporcionant xifratge sense configuració addicional.

%----------------------------------------------------------------------------------------

\section{Compliment normatiu}

El sistema ha estat dissenyat tenint en compte la legislació vigent en matèria de protecció de dades, serveis de la societat de la informació i accessibilitat.

\subsection{Protecció de dades (LOPDGDD)}

El compliment del Reglament General de Protecció de Dades (RGPD) i la Llei Orgànica 3/2018 de Protecció de Dades Personals (LOPDGDD) es garanteix pels següents aspectes:

\begin{itemize}
    \item \textbf{Dades públiques}: Totes les dades utilitzades pel sistema provenen de fonts públiques de la Generalitat de Catalunya, publicades sota llicències obertes. No es tracten dades de caràcter personal.
    \item \textbf{No recollida de dades d'usuaris}: El frontend no recull, emmagatzema ni processa cap informació personal dels usuaris. No hi ha formularis de registre, sistemes de login ni cookies de seguiment.
    \item \textbf{Transparència}: El codi font és completament obert, permetent que qualsevol persona verifiqui què fa el sistema amb les dades.
\end{itemize}

\subsection{LSSICE}

La Llei 34/2002 de Serveis de la Societat de la Informació i de Comerç Electrònic (LSSICE) estableix requisits d'informació per als serveis en línia. En tractar-se d'un projecte acadèmic sense finalitat comercial, els requisits es redueixen, però s'ha inclòs:

\begin{itemize}
    \item Informació clara sobre l'autoria del projecte a la pàgina \texttt{/about}.
    \item Identificació de les fonts de dades utilitzades (APIs Socrata de la Generalitat).
    \item Enllaç al repositori de codi font per a verificació independent.
\end{itemize}

\subsection{Accessibilitat}

L'accessibilitat del sistema s'ha abordat seguint les pautes WCAG 2.1 de forma progressiva:

\begin{itemize}
    \item \textbf{Contrast adequat}: S'han seleccionat colors amb contrast suficient entre text i fons per garantir la llegibilitat.
    \item \textbf{Etiquetes descriptives}: Tots els elements interactius tenen etiquetes clares i comprensibles.
    \item \textbf{Navegació per teclat}: Els elements interactius principals són accessibles mitjançant teclat.
    \item \textbf{Disseny responsiu}: L'aplicació s'adapta a diferents mides de pantalla, facilitant l'accés des de dispositius mòbils.
\end{itemize}

Tot i aquests esforços, es reconeix que algunes visualitzacions complexes (com els gràfics de Recharts) podrien millorar la seva accessibilitat amb text alternatiu més detallat i suport per a lectors de pantalla, aspecte que es podria abordar en futures iteracions.

%----------------------------------------------------------------------------------------

\section{Validació de requisits funcionals}

A continuació es valida el grau d'assoliment de cada requisit funcional definit al Capítol 6:

\begin{itemize}
    \item \textbf{RF-1 (Alta): Extracció automàtica de dades} -- \textbf{ASSOLIT}. Els extractors \texttt{EmbassamentExtractor} i \texttt{PrecipitationExtractor} obtenen les dades més recents de les APIs de Socrata de forma automàtica mitjançant GitHub Actions. El sistema suporta tant el mode incremental com el de sincronització completa.
    
    \item \textbf{RF-2 (Alta): Transformació i neteja de dades} -- \textbf{ASSOLIT}. Els transformadors normalitzen dates a ISO 8601, eliminen duplicats, gestionen valors nuls i estructuren les dades en format JSON estàndard. S'han implementat validacions per detectar i gestionar inconsistències.
    
    \item \textbf{RF-3 (Alta): Emmagatzematge en JSON versionat} -- \textbf{ASSOLIT}. Les dades processades es desen en fitxers JSON amb timestamps al directori \texttt{/etl/data/}. El versionat amb Git proporciona historial complet de canvis, permetent tornar a qualsevol punt anterior.
    
    \item \textbf{RF-4 (Alta): Interfície web amb quatre panells} -- \textbf{ASSOLIT}. S'han implementat els quatre panells de visualització:
    \begin{itemize}
        \item Panell 1 (Mapa): Mapa interactiu amb Leaflet mostrant markers d'embassaments i estacions amb popups informatius i KPIs dinàmics.
        \item Panell 2 (Evolució temporal): Gràfics de línies amb Recharts, selectors de data i funcionalitat de comparació múltiple d'embassaments.
        \item Panell 3 (Correlació): Gràfics de dispersió amb càlcul del coeficient de Pearson i selecció d'estacions meteorològiques.
        \item Panell 4 (Alertes): Taula amb indicadors semàfor basats en percentatge d'ocupació i fletxes de tendència.
    \end{itemize}
    
    \item \textbf{RF-5 (Mitjana): Detecció automàtica de sequera} -- \textbf{ASSOLIT}. S'ha implementat la lògica de classificació basada en llindars: crític ($<$20\%), avís (20-50\%), normal (50-75\%) i òptim ($>$75\%). A més, s'ha afegit un algoritme de regressió lineal sobre les últimes mesures per detectar tendències de millora o deteriorament.
    
    \item \textbf{RF-6 (Mitjana): Automatització amb GitHub Actions} -- \textbf{ASSOLIT}. El workflow \texttt{etl-pipeline.yml} s'executa diàriament a les 02:00 UTC, actualitzant automàticament les dades i desencadenant el desplegament del frontend si cal.
    
    \item \textbf{RF-7 (Baixa): Exportació de dades} -- \textbf{PARCIALMENT ASSOLIT}. Les dades són accessibles directament com a fitxers JSON des del repositori. La funcionalitat d'exportació des del frontend es podria implementar en futures iteracions.
\end{itemize}

%----------------------------------------------------------------------------------------

\section{Validació de requisits no funcionals}

S'avaluen a continuació els requisits no funcionals definits al Capítol 6:

\begin{itemize}
    \item \textbf{RNF-1 (Alta): Rendiment} -- \textbf{ASSOLIT}. L'ETL processa les dades de més de 200 estacions meteorològiques i 9 embassaments en un temps inferior als 5 minuts, complint el requisit. El frontend carrega en menys de 3 segons gràcies a l'optimització de Vite (minificació, tree-shaking) i el carregament asíncron de components amb \texttt{React.lazy}.
    
    \item \textbf{RNF-2 (Alta): Usabilitat} -- \textbf{ASSOLIT}. La interfície s'ha dissenyat seguint principis de simplicitat i claredat, amb navegació intuïtiva, etiquetatge comprensible i elements interactius amb feedback visual immediat. Els indicadors semàfor del Panell 4 proporcionen informació instantània accessible fins i tot per a usuaris no tècnics.
    
    \item \textbf{RNF-3 (Alta): Seguretat} -- \textbf{ASSOLIT}. Totes les connexions a les APIs es realitzen via HTTPS. No s'emmagatzemen credencials al repositori (les APIs són públiques). La configuració externalitzada amb \texttt{.env} permet gestionar paràmetres sensibles si calgués en el futur.
    
    \item \textbf{RNF-4 (Mitjana): Portabilitat} -- \textbf{ASSOLIT}. El frontend és una aplicació estàtica que es pot desplegar a qualsevol servidor web (GitHub Pages, Vercel, Netlify). L'ETL és un script Node.js que pot executar-se en qualsevol entorn amb Node.js 18+, incloent màquines locals, servidors o altres plataformes de CI/CD.
    
    \item \textbf{RNF-5 (Mitjana): Mantenibilitat} -- \textbf{ASSOLIT}. El codi segueix una estructura modular amb separació clara de responsabilitats, convencions de nomenclatura consistents, comentaris explicatius i documentació tècnica. S'utilitzen eines com ESLint i Prettier per garantir la qualitat del codi.
    
    \item \textbf{RNF-6 (Baixa): Escalabilitat} -- \textbf{ASSOLIT}. L'arquitectura desacoblada ETL/frontend permet afegir nous tipus de dades o augmentar el volum sense redisseny complet. La separació mitjançant fitxers JSON com a intermediari facilita aquesta escalabilitat, i el mecanisme de finestra lliscant per a dades de precipitació demostra la capacitat de gestionar grans volums de dades.
\end{itemize}

%----------------------------------------------------------------------------------------

\section{Demostració del producte}

El producte final es troba desplegat i accessible públicament. Es pot consultar directament a l'adreça següent:

\textbf{Aplicació desplegada:} \url{https://alexruizmaso.github.io/opendata-water-visualization}

\textbf{Enllaç al vídeo de demostració:} \url{[enllaç al vídeo]}

%----------------------------------------------------------------------------------------

\section{Valoració global dels resultats}

Els resultats obtinguts amb el desenvolupament d'aquest projecte superen les expectatives inicials en diversos aspectes. El grau d'assoliment dels objectius és elevat: s'ha creat un sistema complet que extreu, processa i visualitza dades hídriques de Catalunya de forma automàtica i accessible.

\subsection{Objectius assolits}

\begin{itemize}
    \item \textbf{Visualització accessible}: S'ha aconseguit transformar dades tècniques complexes en visualitzacions intuïtives que qualsevol ciutadà pot entendre, complint l'objectiu principal de democratitzar l'accés a la informació hídrica.
    \item \textbf{Automatització completa}: El sistema funciona de forma autònoma, actualitzant les dades diàriament sense intervenció manual, gràcies a la integració amb GitHub Actions.
    \item \textbf{Arquitectura robusta}: La separació ETL/frontend mitjançant fitxers JSON ha resultat en un sistema senzill, mantenible i econòmic (cost zero d'infraestructura).
    \item \textbf{Codi obert}: Tot el codi font és públic i reutilitzable, facilitant la contribució comunitària i l'adaptació per a altres regions o tipus de dades.
\end{itemize}

\subsection{Qualitat del producte}

El producte final presenta una qualitat professional, amb una interfície coherent, visualitzacions precises i un rendiment adequat. Les proves unitàries implementades garanteixen el correcte funcionament de cada component, i l'arquitectura modular facilita futures ampliacions.

\subsection{Impacte potencial}

L'impacte potencial del projecte és significatiu en un context de creixent preocupació per la sequera a Catalunya. El sistema proporciona una eina accessible per a:

\begin{itemize}
    \item \textbf{Ciutadans}: Per entendre l'estat dels recursos hídrics i prendre consciència sobre la situació de sequera.
    \item \textbf{Professionals i investigadors}: Per analitzar correlacions entre precipitació i nivells d'embassaments amb dades actualitzades.
    \item \textbf{Administracions}: Per disposar d'una eina de referència accessible que complementa els sistemes oficials de monitoratge.
\end{itemize}

\subsection{Limitacions}

Es reconeixen les següents limitacions:

\begin{itemize}
    \item La dependència de GitHub Actions limita l'execució de l'ETL a una freqüència mínima d'aproximadament 5 minuts per execució, suficient per a les necessitats actuals però potencialment insuficient per a monitoratge en temps real.
    \item L'absència d'una base de dades real limita la capacitat de consultes complexes sobre l'historial complet de dades.
    \item Algunes funcionalitats avançades com l'exportació de dades des del frontend s'han deixat per a futures iteracions.
\end{itemize}

Malgrat aquestes limitacions, el sistema compleix els seus objectius principals i proporciona una base sòlida per a futures millores.

%----------------------------------------------------------------------------------------